% =====================================================================
%  Rozdział 3 — Projekt i wykonanie prototypu (budżet: 6,5 str., ok. 1500 słów)
%  Plan akapitów: wspolne/rozdzial3_plan.md (oznaczenia T1–T20 = brakujące dane)
%  Wersja robocza 09.11.2026: napisane 3.1 i 3.2 (budowa struktury, szkło 1 mm),
%  pozostałe podrozdziały jako nagłówki z \TODO.
% =====================================================================
\chapter{Projekt i~wykonanie prototypu}
\label{ch:projekt}

W~rozdziale opisano projekt struktury krzemowo-szklanej źródła elektronów
oraz kolejne etapy jej wytworzenia, aż do zmontowanego prototypu. Zasady
zastosowanych procesów technologicznych omówiono w~\pref{sec:technologia},
a~stanowisko pomiarowe w~\rozdzref{ch:stanowisko}.

% ---------------------------------------------------------------------
\section{Założenia projektowe i~koncepcja struktury}
\label{sec:koncepcja}

Prototyp miał spełniać następujące wymagania:
\begin{itemize}[nosep]
  \item wytworzenie wiązki elektronów przez emisję polową z~nanorurek
        węglowych przy napięciu ekstrakcyjnym $\UB$ rzędu \SI{1}{\kilo\volt},
  \item przyspieszenie elektronów do energii wystarczającej do wzbudzenia
        linii charakterystycznych pierwiastków lekkich (w~badaniach
        $\UT$ do \SI{4}{\kilo\volt}, docelowo \SIrange{10}{18}{\kilo\volt}),
  \item izolacja elektryczna elektrod przy tych napięciach,
  \item wykonanie wyłącznie procesami MEMS: fotolitografią, trawieniem
        krzemu i~bondingiem anodowym krzem--szkło,
  \item możliwość późniejszej hermetyzacji struktury.
\end{itemize}

Strukturę zaprojektowano jako stos płaskich elementów krzemowych
rozdzielonych szkłem borokrzemowym (\rysref{fig:przekroj-struktury}).
Licząc od dołu, składa się ona z~emitera (E) z~ostrzem uformowanym
z~nanorurek węglowych, szkła dystansowego o~grubości \SI{1}{\milli\metre},
bramki (elektrody ekstrakcyjnej, B), drugiego szkła dystansowego oraz
targetu (T). Każda elektroda krzemowa jest osobnym elementem
z~własnym kontaktem elektrycznym, a~szkło oddziela je galwanicznie
i~ustala odległości między nimi. Wiązkę ogniskuje magnes trwały NdFeB
umieszczony przy emiterze
\TODO{T1: położenie magnesu i~ewentualne dodatkowe elementy wg szkicu}.

\begin{figure}[htbp]
  \centering
  \begin{tikzpicture}[x=1cm, y=1cm, font=\small]
    % kolory
    \colorlet{si}{gray!35}
    \colorlet{gl}{cyan!12}
    % target
    \fill[si] (0,4.4) rectangle (8,4.9);
    % szkło 2 z otworem
    \fill[gl] (0,3.4) rectangle (3,4.4);
    \fill[gl] (5,3.4) rectangle (8,4.4);
    % bramka z aperturą
    \fill[si] (0,2.9) rectangle (3.7,3.4);
    \fill[si] (4.3,2.9) rectangle (8,3.4);
    % szkło 1 z otworem
    \fill[gl] (0,1.9) rectangle (3,2.9);
    \fill[gl] (5,1.9) rectangle (8,2.9);
    % emiter
    \fill[si] (0,1.4) rectangle (8,1.9);
    % wieża (ostrze) CNT
    \fill[black!75] (3.7,1.9) -- (4,2.45) -- (4.3,1.9) -- cycle;
    % magnes (poglądowo, położenie do potwierdzenia)
    \fill[orange!40] (3.2,0.7) rectangle (4.8,1.4);
    % obrysy
    \foreach \y in {1.4,1.9,2.9,3.4,4.4,4.9} {\draw[gray] (0,\y) -- (8,\y);}
    \draw (0,1.4) rectangle (8,4.9);
    \draw (3.2,0.7) rectangle (4.8,1.4);
    % wiązka
    \draw[-{Stealth[length=2.5mm]}, blue!70!black, thick, dashed]
      (4,2.5) -- (4,4.38);
    % opisy
    \foreach \y/\t in {4.65/target (T),
                       3.9/szkło borokrzemowe,
                       3.15/bramka (B),
                       2.4/szkło borokrzemowe \SI{1}{\milli\metre},
                       1.65/emiter (E),
                       1.05/magnes NdFeB}
      {\node[anchor=west] at (8.3,\y) {\t};}
    \node[anchor=east] at (3.4,2.3) {ostrze CNT};
    \draw[gray] (3.42,2.25) -- (3.82,2.1);    \node[anchor=west, blue!70!black] at (4.1,3.75) {$e^-$};
    % wymiar
    \draw[{Stealth[length=2mm]}-{Stealth[length=2mm]}] (-0.3,1.9) -- (-0.3,2.9)
      node[midway, left] {$d$};
  \end{tikzpicture}
  \caption{Przekrój struktury prototypu (schemat poglądowy, bez zachowania
    skali): E -- emiter z~ostrzem CNT, B -- bramka (elektroda ekstrakcyjna),
    T -- target; $d$ -- odległość emiter--bramka \opracowanieWlasne}
  \label{fig:przekroj-struktury}
\end{figure}
% TODO(rys.): podmienić na wersję zgodną ze szkicem kris (T1); sprawdzić
% grubość drugiego szkła, kształt emitera, wymiary otworów i położenie magnesu.

W~porównaniu ze źródłem promieniowania rentgenowskiego opisanym
w~\cite{urbanski2026phd} konstrukcja jest uproszczona: nie ma
elektrody ogniskującej ani zintegrowanej mikropompy, a~wiązka pada
bezpośrednio na target, który pełni rolę nośnika badanej próbki
(lub warstwy luminoforu), a~nie targetu transmisyjnego.
Sposób zasilania elektrod opisano w~\rozdzref{ch:stanowisko}.
\TODO{T2: które elementy przejęto z~projektu źródła RTG, a~które zaprojektowano od nowa}

% ---------------------------------------------------------------------
\section{Dobór materiałów i~grubości elementów}
\label{sec:materialy}

Elektrody wykonano z~płytek krzemu monokrystalicznego o~orientacji
(100), ponieważ kształt emitera i~otworów uzyskuje się trawieniem
anizotropowym (\pref{subsec:koh}), a~krzem jest jednocześnie materiałem
konstrukcyjnym i~przewodzącym.
\TODO{T4: grubość, typ i~rezystywność płytek Si}

Jako dystansownik między emiterem a~bramką zastosowano szkło
borokrzemowe o~grubości \SI{1}{\milli\metre}. Grubość szkła wyznacza
bezpośrednio odległość emiter--bramka $d$, a~więc natężenie pola
elektrycznego przy emiterze. Przy $d \approx \SI{1}{\milli\metre}$
i~$\UB = \SI{1}{\kilo\volt}$ średnie natężenie pola wynosi ok.
\SI{1}{\volt\per\micro\metre}; lokalne pole na wierzchołkach nanorurek
jest wielokrotnie większe dzięki współczynnikowi wzmocnienia pola
(\pref{sec:emisja}). Ostrze CNT o~wysokości $h$ skraca odległość
wierzchołka od bramki do $d - h$, co dodatkowo zwiększa pole przy
wierzchołku.
\TODO{T7: wysokość ostrza CNT $h$ i~średnica apertury bramki}

Drugim zadaniem szkła jest izolacja elektrod. W~próżni ograniczeniem
jest przede wszystkim przeskok wzdłuż powierzchni szkła między krawędziami
sąsiednich elektrod, dlatego grubość dystansownika wyznacza też drogę
upływu. Między bramką a~targetem występuje różnica potencjałów
$\UT - \UB$, a~między emiterem a~bramką $\UB$. Przy napięciach
stosowanych w~badaniach (do \SI{4}{\kilo\volt}) szkło o~grubości
\SI{1}{\milli\metre} zapewniało wystarczającą izolację. Przy docelowych
napięciach \SIrange{10}{18}{\kilo\volt} potrzebne są dodatkowe
zabezpieczenia, np. zalanie krawędzi struktury żywicą epoksydową, jak
w~\cite{urbanski2026phd}.
\TODO{T5: typ szkła (Borofloat~33?), grubość drugiego szkła, powód wyboru 1~mm}

Szkło borokrzemowe ma współczynnik rozszerzalności cieplnej zbliżony do
krzemu i~zawiera jony sodu, dzięki czemu można je łączyć z~krzemem
bondingiem anodowym (\pref{subsec:bonding}). Duża odległość między
emiterem a~targetem sprzyja jednak rozbieganiu się wiązki, co uzasadnia
zastosowanie magnesu ogniskującego (\pref{sec:wiazka}).
Wymiary wszystkich elementów zestawiono w~\tabref{tab:elementy-struktury}.

\begin{table}[htbp]
  \centering
  \caption{Elementy struktury prototypu (od dołu)}
  \label{tab:elementy-struktury}
  \begin{tabular}{@{}llll@{}}
    \toprule
    Element & Materiał & Grubość & Funkcja \\
    \midrule
    emiter (E)        & Si (100) + ostrze CNT & \TODO{T4}          & źródło elektronów \\
    dystansownik      & szkło borokrzemowe   & \SI{1}{\milli\metre} & izolacja, odległość $d$ \\
    bramka (B)        & Si (100)             & \TODO{T4}           & elektroda ekstrakcyjna \\
    dystansownik      & szkło borokrzemowe   & \TODO{T5}           & izolacja B--T \\
    target (T)        & \TODO{T8}            & \TODO{T8}           & nośnik próbki / luminoforu \\
    magnes ogniskujący & NdFeB               & --                  & ogniskowanie wiązki \\
    \bottomrule
  \end{tabular}
\end{table}

% ---------------------------------------------------------------------
\section{Projekt emitera, bramki i~targetu}
\label{sec:projekt-elektrod}
\TODO{T6--T8: wymiary emitera, apertura bramki, target; magnes jakościowo (plan p.~3.3)}

\section{Projekt masek fotolitograficznych}
\label{sec:maski}
\TODO{T9--T10: liczba i~przeznaczenie masek, reguły projektowe, rozkład na płytce (plan p.~3.4)}

\section{Wytwarzanie elementów krzemowych}
\label{sec:wytwarzanie-si}
\TODO{T11--T14: warstwa maskująca, fotolitografia, trawienie KOH, kontrola wymiarów (plan p.~3.5)}

\section{Wykonanie elementów szklanych}
\label{sec:szklo}

Dystansowniki wykonano z~płytki szkła borokrzemowego o~grubości
\SI{1}{\milli\metre}. Najpierw wytrawiono w~niej otwory o~średnicy
\SI{0.6}{\milli\metre}, przez które przechodzi wiązka elektronów,
a~następnie płytkę pocięto piłą na pojedyncze elementy. Taka kolejność
pozwala obrabiać całą płytkę naraz i~zachować jednakowe położenie
otworów we wszystkich dystansownikach. Przed bondingiem elementy
oczyszczono, ponieważ zanieczyszczenia i~nierówności powierzchni
uniemożliwiają szczelne połączenie (\pref{subsec:bonding}).
\TODO{T15: roztwór i~maska do trawienia otworów w~szkle, sposób czyszczenia}

\section{Formowanie emitera CNT}
\label{sec:nanoszenie-cnt}

Emiter ma postać ostrza (wieży) z~nanorurek węglowych na płaskim
elemencie krzemowym. Zawiesinę nanorurek nakraplano pipetą na element
pomocniczy, który przykładano do emitera nagrzanego na płycie grzejnej
do \SI{240}{\celsius}. Element pomocniczy odciągano następnie szybkim
ruchem od powierzchni emitera. Zawiesina rozciągała się między oboma
elementami i~zasychała w~kontakcie z~gorącym podłożem, tworząc wieżę
nanorurek wydłużoną w~kierunku przyszłej bramki. Kształt ostrza
zwiększa lokalne natężenie pola na jego wierzchołku i~skraca odległość
od bramki (\pref{sec:materialy}).
\TODO{T16: skład zawiesiny (rodzaj CNT, rozpuszczalnik, PVP?), materiał elementu pomocniczego, wysokość i~średnica uzyskanej wieży}
\TODO{T17: czy emiter formowano przed bondingiem, czy po nim}

\section{Bonding anodowy i~montaż prototypu}
\label{sec:bonding-montaz}

Elementy połączono bondingiem anodowym w~atmosferze. Stół grzejny
nagrzewano do \SI{450}{\celsius}, czyli do górnej granicy typowego
zakresu temperatur procesu (\pref{subsec:bonding}). Stół był
jednocześnie elektrodą uziemiającą, a~kolejne warstwy stosu znajdowały
się coraz dalej od niego, więc ich temperatura była niższa niż
temperatura stołu. Wysoka temperatura stołu zapewniała, że także
najwyżej położone połączenia osiągały temperaturę wystarczającą do
przewodnictwa jonowego szkła. Między elementy wprowadzono dodatkowo
przekładkę, która zmniejszała dystans między łączonymi elementami.
\TODO{T18: kolejność łączenia warstw, sposób justowania otworów względem ostrza; czym była przekładka i~gdzie leżała}
\TODO{T19: napięcie i~czas bondingu, czy wszystkie połączenia były trwałe}

\TODO{T20: montaż -- położenie magnesu, kontakty elektryczne do E, B, T, luminofor i~próbka na targecie, wymiary i~masa prototypu}

Wszystkie połączenia wykonano w~atmosferze, dlatego wnętrze struktury
nie jest zamknięte w~próżni, a~prototyp jest niehermetyczny. Badania
prowadzono w~komorze próżniowej przy ciśnieniu ok.
\SI{1e-5}{\hecto\pascal} (\rozdzref{ch:stanowisko}). Rozwiązaniem
docelowym jest wykonanie ostatniego połączenia bondingiem anodowym
w~próżni i~podtrzymywanie próżni mikropompą; trudności tego procesu
omówiono w~\pref{subsec:bonding}.
